\section{Evaluation}
\label{sec:evaluation}
We evaluate whether token scores recover the evidence used in a response, predict the effect of removing that evidence, and can be computed efficiently.

\paragraph{Tasks and methods.}
We use Qwen3-8B and the 1,243 examples released with FlashTrace \citep{pan2026flashtrace}: ten RULER tasks with 100 examples each, HotpotQA with 48, MATH with 100, and MoreHopQA with 95.\footnote{\url{https://github.com/wbopan/flashtrace/releases/tag/table1-data-v1}} Needle-in-a-haystack (NIAH) tasks retrieve information among distractors. Multi-query tasks (MQ-Q2/Q4/Q8) ask for 2, 4, or 8 keys; multi-value tasks (MV-V2/V4/V8) ask for 2, 4, or 8 values of one key. Variable Tracking (VT) follows variable assignments: H is the number of binding hops and C the number of parallel chains. HotpotQA and MoreHopQA require combining facts, while MATH contains mathematical reasoning problems. The complete stored reasoning and answer plus EOS is scored for faithfulness. We compare \dt{} (DT) with FlashTrace (FT), Perturbation, REAGENT, Context Length Probing (CLP), Information Flow Routes (IFR), and Attention-aware Layer-wise Relevance Propagation (AttnLRP). RISE/MAS and NIAH baseline results use the published aggregates; VT/HotpotQA recovery uses the seven-method re-evaluation in Appendix~\ref{app:evaluation}.

\paragraph{Evidence recovery.}
Recall is the fraction of annotated evidence recovered among the highest-ranked sources. A token budget limits how much input the ranking may select. NIAH uses the top 10\% of eligible input tokens, the source positions marked for evaluation in the released examples. VT measures token Recall in the passage body, with budgets of 10\%, 10\%, 20\%, and 30\% for H2, H4, H6, and H10. HotpotQA measures Recall of supporting sentences selected within a 10\% body-token budget. VT recovery attributes the answer tokens; HotpotQA attributes the full response and ranks sentences by their summed signed scores. FT uses three recursive tracing steps (K3) for VT/HotpotQA recovery. Table~\ref{tab:recovery} gives these values; the first-page radar shows the same Recall values with each axis's minimum and maximum labeled.

\paragraph{Deletion faithfulness.}
The released deletion test masks high-ranked source tokens with a baseline token in 20 steps, rescoring the same complete response after each step. The score is normalized between the full-input and fully-deleted endpoints. A good ranking causes an early score drop because it removes evidence the model uses. The RISE deletion score summarizes the area under this curve. Magnitude Aligned Scoring (MAS) also measures how the removed attribution mass aligns with the score change \citep{pan2026flashtrace}. Lower values are better for both metrics. DT uses signed scores for RISE and their positive part for MAS. Appendix~\ref{app:evaluation} specifies the ranking and normalization.

\begin{table}[t]
\centering\small
\caption{RISE $\downarrow$ on the complete released Qwen3-8B tasks. FT is FlashTrace \citep{pan2026flashtrace} and DT is \dt{}. DT uses signed ranking. FT uses the same evaluation examples as DT. $n$ is the released task size. Bold marks the lowest observed mean across all seven methods.}
\label{tab:full}
\setlength{\tabcolsep}{3.0pt}
\begin{tabular}{@{}lr rrrrrrr@{}}
\toprule
Task & $n$ & Perturbation & REAGENT & CLP & IFR & AttnLRP & FT & DT \\
\midrule
MQ-Q2 & 100 & 0.0945 & 0.1174 & 0.0977 & 0.0747 & 0.1955 & \textbf{0.0682} & 0.0723 \\
MQ-Q4 & 100 & 0.2389 & 0.2599 & 0.2534 & 0.1147 & 0.2628 & \textbf{0.1134} & 0.1448 \\
MQ-Q8 & 100 & 0.4986 & 0.4865 & 0.5096 & 0.3709 & 0.3766 & \textbf{0.3519} & 0.4286 \\
MV-V2 & 100 & 0.1343 & 0.1876 & 0.1559 & 0.0691 & 0.1403 & 0.0686 & \textbf{0.0597} \\
MV-V4 & 100 & 0.1862 & 0.2107 & 0.2166 & 0.0726 & 0.1934 & 0.0704 & \textbf{0.0587} \\
MV-V8 & 100 & 0.3514 & 0.3685 & 0.3929 & 0.2051 & 0.2852 & 0.1828 & \textbf{0.1159} \\
\midrule
VT-H2-C3 & 100 & 0.3844 & 0.4378 & 0.3742 & 0.1609 & 0.3191 & 0.1317 & \textbf{0.1003} \\
VT-H4-C1 & 100 & 0.3538 & 0.3971 & 0.3277 & 0.1018 & 0.3239 & 0.1102 & \textbf{0.1013} \\
VT-H6-C1 & 100 & 0.4584 & 0.4860 & 0.4225 & 0.1245 & 0.3378 & 0.1224 & \textbf{0.1109} \\
VT-H10-C1 & 100 & 0.4663 & 0.4953 & 0.4510 & 0.1526 & 0.3573 & 0.1427 & \textbf{0.1403} \\
\midrule
HotpotQA & 48 & 0.1333 & 0.1445 & 0.1013 & 0.0745 & 0.1548 & 0.0331 & \textbf{0.0275} \\
MATH & 100 & 0.3802 & 0.3878 & 0.3616 & 0.3537 & 0.3684 & 0.3487 & \textbf{0.2321} \\
MoreHopQA & 95 & 0.2489 & 0.2652 & 0.2487 & 0.1455 & 0.2857 & 0.1281 & \textbf{0.0703} \\
\bottomrule
\end{tabular}
\end{table}

\begin{table}[t]
\centering\small
\caption{MAS $\downarrow$ on the same released tasks and seven methods as Table~\ref{tab:full}. DT uses its positive source contributions. Bold marks the lowest observed mean.}
\label{tab:mas}
\setlength{\tabcolsep}{3.0pt}
\begin{tabular}{@{}lr rrrrrrr@{}}
\toprule
Task & $n$ & Perturbation & REAGENT & CLP & IFR & AttnLRP & FT & DT \\
\midrule
MQ-Q2 & 100 & 0.1442 & 0.1966 & 0.1663 & 0.1397 & 0.3261 & 0.1628 & \textbf{0.1176} \\
MQ-Q4 & 100 & 0.3270 & 0.3574 & 0.3203 & \textbf{0.1770} & 0.4509 & 0.1937 & 0.2011 \\
MQ-Q8 & 100 & 0.7093 & 0.6942 & 0.6566 & 0.4600 & 0.6020 & \textbf{0.4271} & 0.6047 \\
MV-V2 & 100 & 0.1870 & 0.2756 & 0.2159 & 0.1343 & 0.2290 & 0.1569 & \textbf{0.1046} \\
MV-V4 & 100 & 0.2436 & 0.2905 & 0.2804 & 0.1416 & 0.3248 & 0.1605 & \textbf{0.1058} \\
MV-V8 & 100 & 0.4580 & 0.4938 & 0.5114 & 0.2746 & 0.4752 & 0.2477 & \textbf{0.1717} \\
\midrule
VT-H2-C3 & 100 & 0.5513 & 0.6680 & 0.4900 & 0.2309 & 0.5214 & 0.1873 & \textbf{0.1486} \\
VT-H4-C1 & 100 & 0.5168 & 0.6031 & 0.4199 & 0.1479 & 0.5475 & 0.1656 & \textbf{0.1406} \\
VT-H6-C1 & 100 & 0.6837 & 0.7453 & 0.5648 & 0.1727 & 0.5723 & 0.1787 & \textbf{0.1537} \\
VT-H10-C1 & 100 & 0.7013 & 0.7405 & 0.5967 & 0.2011 & 0.5915 & 0.1983 & \textbf{0.1842} \\
\midrule
HotpotQA & 48 & 0.2203 & 0.2351 & 0.1895 & 0.1655 & 0.2493 & 0.1276 & \textbf{0.1002} \\
MATH & 100 & 0.5202 & 0.5358 & 0.4907 & 0.4897 & 0.5244 & 0.4461 & \textbf{0.3323} \\
MoreHopQA & 95 & 0.3470 & 0.3679 & 0.3478 & 0.2277 & 0.4579 & 0.2049 & \textbf{0.1546} \\
\bottomrule
\end{tabular}
\end{table}

\begin{table}[t]
\centering\footnotesize
\caption{Evidence recovery ($\%$, higher is better) at the listed token budgets. NIAH/VT report token Recall. HotpotQA reports Full-support Recall@30\%, the fraction of examples for which every official supporting fact is selected. Sentence scores sum positive token contributions. All body tokens are included in the budget calculation. VT macro averages its four rows. FT uses K3 for recovery. Bold marks the highest observed mean.}
\label{tab:recovery}
\setlength{\tabcolsep}{2.2pt}
\begin{tabular}{@{}lrr rrrrrrr@{}}
\toprule
Task & $n$ & Budget & Perturbation & REAGENT & CLP & IFR & AttnLRP & FT & DT \\
\midrule
MQ-Q2 & 100 & 10\% & 39.05 & 24.35 & 39.85 & 47.10 & 21.49 & 46.32 & \textbf{67.24} \\
MQ-Q4 & 100 & 10\% & 9.04 & 8.49 & 8.58 & 32.84 & 20.42 & 39.58 & \textbf{47.16} \\
MQ-Q8 & 100 & 10\% & 0.97 & 0.47 & 0.84 & 1.19 & 7.63 & 7.72 & \textbf{13.59} \\
MV-V2 & 100 & 10\% & 25.48 & 18.02 & 20.67 & 57.50 & 25.44 & 53.46 & \textbf{77.79} \\
MV-V4 & 100 & 10\% & 16.05 & 15.60 & 14.61 & 45.24 & 24.25 & 49.37 & \textbf{56.79} \\
MV-V8 & 100 & 10\% & 8.02 & 7.37 & 7.29 & 17.87 & 15.88 & 19.52 & \textbf{25.19} \\
\midrule
VT-H2-C3 & 100 & 10\% & 26.69 & 26.66 & 29.10 & 81.16 & 83.37 & 72.41 & \textbf{93.81} \\
VT-H4-C1 & 100 & 10\% & 18.86 & 17.58 & 20.53 & 70.09 & 70.54 & 68.32 & \textbf{71.31} \\
VT-H6-C1 & 100 & 20\% & 33.57 & 32.47 & 33.85 & 83.74 & 84.57 & 82.15 & \textbf{95.37} \\
VT-H10-C1 & 100 & 30\% & 40.70 & 42.48 & 44.04 & 82.66 & 77.33 & 83.41 & \textbf{90.51} \\
VT macro & 400 & mixed & 29.96 & 29.80 & 31.88 & 79.41 & 78.95 & 76.57 & \textbf{87.75} \\
\midrule
HotpotQA & 48 & 30\% & 35.42 & 35.42 & 64.58 & 85.42 & \textbf{87.50} & 75.00 & 85.42 \\
\bottomrule
\end{tabular}
\end{table}


\paragraph{Attribution quality.}
DT has the lowest RISE on 9 of 13 tasks and the lowest MAS on 11 (Tables~\ref{tab:full} and~\ref{tab:mas}). Against FT, the corresponding improvements cover 10 and 12 tasks. The largest absolute RISE reduction is on MATH, from 0.3484 to 0.2029. DT achieves the highest Recall on all six NIAH tasks, exceeding FT by 3.96--19.60 percentage points. At the stated VT budgets, its macro Recall is 84.18\%, compared with 76.57\% for FT K3. On HotpotQA, IFR reaches 75.87\% and DT reaches 72.05\% (Table~\ref{tab:recovery}).

\begin{table}[t]
\centering\small
\caption{Attribution quality on Qwen3.5-9B. All tasks use the full set of released examples. Recovery is token Recall for NIAH/VT and Full-support Recall@30\% for HotpotQA, as defined in Table~\ref{tab:recovery}. FT denotes FlashTrace, with K3 for Recall and K1 for RISE/MAS. DT uses signed RISE and positive-part MAS. Bold marks the better mean. Dashes indicate undefined recovery.}
\label{tab:qwen35-quality}
\setlength{\tabcolsep}{4.0pt}
\begin{tabular}{@{}lrr rr rr rr@{}}
\toprule
Task & $n$ & Budget & \multicolumn{2}{c}{Recall (\%) $\uparrow$} & \multicolumn{2}{c}{RISE $\downarrow$} & \multicolumn{2}{c}{MAS $\downarrow$} \\
\cmidrule(lr){4-5}\cmidrule(lr){6-7}\cmidrule(l){8-9}
 & & & FT & DT & FT & DT & FT & DT \\
\midrule
MQ-Q2 & 100 & 10\% & 52.92 & \textbf{71.94} & \textbf{0.0982} & 0.0996 & 0.2549 & \textbf{0.1306} \\
MQ-Q4 & 100 & 10\% & 38.56 & \textbf{54.18} & \textbf{0.1414} & 0.1579 & 0.2651 & \textbf{0.2242} \\
MQ-Q8 & 100 & 10\% & 10.72 & \textbf{14.63} & \textbf{0.3704} & 0.4028 & \textbf{0.4390} & 0.5891 \\
MV-V2 & 100 & 10\% & 60.51 & \textbf{83.53} & 0.0974 & \textbf{0.0804} & 0.2467 & \textbf{0.1080} \\
MV-V4 & 100 & 10\% & 46.19 & \textbf{61.42} & 0.1284 & \textbf{0.1130} & 0.2538 & \textbf{0.1510} \\
MV-V8 & 100 & 10\% & 22.27 & \textbf{26.36} & 0.2508 & \textbf{0.2206} & 0.3264 & \textbf{0.2992} \\
\midrule
MATH & 100 & --- & --- & --- & 0.3777 & \textbf{0.2638} & 0.4837 & \textbf{0.3368} \\
MoreHopQA & 95 & --- & --- & --- & 0.2093 & \textbf{0.1429} & 0.2894 & \textbf{0.2001} \\
\midrule
VT H2-C3 & 100 & 10\% & 67.80 & \textbf{94.20} & 0.1514 & \textbf{0.1095} & 0.2224 & \textbf{0.1295} \\
VT H4-C1 & 100 & 10\% & 69.07 & \textbf{71.89} & 0.1161 & \textbf{0.1013} & 0.1888 & \textbf{0.1343} \\
VT H6-C1 & 100 & 20\% & 83.06 & \textbf{97.35} & 0.1300 & \textbf{0.1154} & 0.2051 & \textbf{0.1504} \\
VT H10-C1 & 100 & 30\% & 81.90 & \textbf{94.23} & \textbf{0.1537} & 0.1561 & 0.2307 & \textbf{0.2040} \\
HotpotQA & 48 & 30\% & 58.33 & \textbf{95.83} & 0.2154 & \textbf{0.1470} & 0.3125 & \textbf{0.2367} \\
\bottomrule
\end{tabular}
\end{table}


\paragraph{Qwen3.5-9B results.}
Table~\ref{tab:qwen35-quality} compares DT and FT with the DT rules identified in each block. The symmetric-GDN block contains eight tasks and 72 examples; DT achieves higher Recall on all six NIAH tasks, lower RISE on six of eight tasks, and lower MAS on seven. On MATH and MoreHopQA, DT lowers both RISE and MAS. The original-GDN block adds 400 VT examples and 48 HotpotQA examples. DT lowers MAS on all five tasks. On HotpotQA, its Recall is 76.56\% versus FT's 57.12\%, with lower RISE and MAS.

\paragraph{Attribution time and memory.}
The first-page timing panel varies response length on Qwen3-8B. DT, FT, and multi-hop FT (FT-mh) are measured on one MetaX C550; the other curves show the published timing references. Crosses mark out-of-memory runs. On the 16-example Qwen3-8B benchmark, the sum of per-example median attribution times is 11.40 seconds for DT and 12.16 seconds for one-hop FT; peak allocated memory is 18.76 and 21.49 GB. On Qwen3.5-9B, processing the 16 examples in batches of two with GPU checkpoints reduces time from 16.68 seconds for serial execution to 12.33 seconds, a 35.3\% throughput increase. Peak memory is 21.64 and 23.88 GB, respectively. Times cover complete attribution calls; Appendix~\ref{app:cost} gives the measurement protocol.

\paragraph{Reading the token scores.}
Figures~\ref{fig:retrieval} and~\ref{fig:multihop} show retrieval and multi-hop examples with each model's own scores. Teal marks positive contributions and coral negative contributions to the complete response score. The retrieval example links source keys to their numbers and the query naming them. The multi-hop example follows a voice actor to a film and its director, whose surname determines the answer. Each deletion curve compares DT with one-hop FT on the same model and response. These views connect token-level scores to the evidence selected and the effect of deleting it.
